\section{Ancaman Keamanan}

Ancaman keamanan terhadap sistem kriptografi dapat diklasifikasikan berdasarkan jenis serangan:

\begin{table}[!htbp]
\centering
\caption{Jenis serangan kriptografi berdasarkan informasi yang dimiliki penyerang}
\label{tab:jenis-serangan}
\begin{tabular}{@{}>{\raggedright\arraybackslash}p{0.22\textwidth}>{\raggedright\arraybackslash}p{0.42\textwidth}>{\raggedright\arraybackslash}p{0.30\textwidth}@{}}
\toprule
\textbf{Jenis Serangan} & \textbf{Informasi Penyerang} & \textbf{Tingkat Kesulitan} \\
\midrule
Ciphertext-only & Hanya ciphertext & Paling sulit bagi penyerang \\
Known-plaintext & Pasangan plaintext--ciphertext & Sedang \\
Chosen-plaintext & Dapat memilih plaintext, dapat ciphertext & Lebih mudah \\
Chosen-ciphertext & Dapat memilih ciphertext, dapat plaintext & Paling mudah bagi penyerang \\
\bottomrule
\end{tabular}
\end{table}

\begin{enumerate}
\item \textbf{Serangan ciphertext-only:} Penyerang hanya memiliki akses ke ciphertext. Ini adalah skenario terlemah bagi penyerang.
\item \textbf{Serangan known-plaintext:} Penyerang memiliki pasangan plaintext--ciphertext dan mencoba memperoleh kunci.
\item \textbf{Serangan chosen-plaintext:} Penyerang dapat memilih plaintext dan memperoleh ciphertext yang sesuai.
\item \textbf{Serangan chosen-ciphertext:} Penyerang dapat memilih ciphertext dan memperoleh plaintext yang sesuai.
\end{enumerate}

Selain itu, ancaman juga mencakup \textbf{man-in-the-middle} (penyerang menyisipkan diri antara Alice dan Bob), \textbf{replay attack} (mengulang pesan yang pernah dikirim), dan \textbf{denial of service} (membuat layanan tidak tersedia) \cite{menezes1996,schneier1996}.
